ただし，本例の数値は架空であり，実測に基づく計算量の検証ではない．実測レポートでは，入力の並び，ピボット選択，処理系，CPU負荷，メモリ階層などを記録し，反復回数と不確かさを踏まえて評価する必要がある．また，4点の実行時間だけから漸近計算量を証明することはできないため，本例では「整合した」と表現し，因果関係を断定しなかった．
